% Supplementary Materials, built on its own:
%
%   "D:/environment/tools/tectonic/tectonic.exe" supplementary.tex   (run from paper/)
%
% Same class and same bibliography as main.tex; the only difference is the `supfile'
% article type, which mdpi.cls uses to prefix every table and figure number with S and to
% set the S-page header.  Every item here was in the main manuscript until the sixth
% revision pass and moved out under the referee's length table (report section G); each
% carries the source comment it had there, and nothing here was recomputed.
\documentclass[mathematics,supfile,submit,moreauthors]{Definitions/mdpi}
\usepackage{xurl}

% Match the lowercase theorem environments used by the included theory section.
\let\proposition\Proposition
\let\endproposition\endProposition
\let\lemma\Lemma
\let\endlemma\endLemma

% EPS logo under XeTeX: same substitution as main.tex; see the comment there.
\makeatletter
\let\mdpi@includegraphics\includegraphics
\renewcommand{\includegraphics}[2][]{%
  \IfEndWith{#2}{.eps}%
    {\StrGobbleRight{#2}{4}[\mdpi@gfile]%
     \mdpi@includegraphics[#1]{\mdpi@gfile.pdf}}%
    {\mdpi@includegraphics[#1]{#2}}%
}
\makeatother

% The same marker the manuscript uses for a development figure.
\newcommand{\devnum}[1]{#1}
% And the one for a figure from the sealed terms.
\newcommand{\sealednum}[1]{#1}

% mdpi.cls prefixes tables and figures with S under the `supfile' article type but leaves
% sections plain; the manuscript's back matter and its cross-references call them S1--S3.
\renewcommand{\thesection}{S\arabic{section}}

%=================================================================
\firstpage{1}
\makeatletter
\setcounter{page}{\@firstpage}
\makeatother
\pubvolume{1}
\issuenum{1}
\articlenumber{0}
\pubyear{2026}
\copyrightyear{2026}

\Title{Prediction-Driven Non-Preemptive Scheduling with Bounded Overtaking for
Shared Execution Services under Deadline-Driven Bursty Load}

% The author block is the manuscript's and is filled in there; it is repeated here
% unchanged so that the two files stay in step.
\Author{[First Author Name] $^{1}$ and [Second Author Name] $^{2,}$*}
\AuthorNames{[First Author Name] and [Second Author Name]}

\address{%
$^{1}$ \quad [First author: department, institution, city and postcode, country]; [first author e-mail]\\
$^{2}$ \quad [Second author: department, institution, city and postcode, country]; [second author e-mail]}

\corres{Correspondence: [corresponding author e-mail]}

%=================================================================
\begin{document}

\maketitle

This file holds supporting material for the manuscript: the visibility audit of the feature pipeline and the predictor
scores on the four non-judging traces (Section~S1); the field audit of the main judging
trace and two checks on the continuous-integration refit (Section~S2); five tables of
experimental detail, the selected budget points, the two budget shapes the main guard
table no longer prints, the identity residual on all nine cells and the single-server
configuration (Section~S3); and two deferred proofs, of the workload comparison and of
the consistency proposition (Section~S4). Section numbers here carry an S prefix and every
table is numbered S. A cross-reference without an S --- ``Section~8.2'', ``Table~8'',
``Theorem~1'' --- points into the main manuscript, since a \verb|\ref| across the two
files does not resolve. The traces are the ones the manuscript's Table~3 names: the
CodeBench judging logs \cite{codebench_dataset}, the Azure Functions invocation trace
\cite{azure2021trace}, the Intel Netbatch logs \cite{netbatch2012trace} and the two
Firefox continuous-integration pools collected from Mozilla's public APIs
\cite{taskcluster_docs,treeherder_docs}. The additional reproducibility specification is Section~S5; capacity, physical
queueing, feedback and LPC-EGEE results are Section~S6. Additional theory is Section~S7. Every empirical figure quoted here is a development
figure in the sense of Section~8.2 of the manuscript: nothing in this file comes from a
sealed part of any trace, and nothing here was recomputed for it.

%=================================================================
\section{Cost Prediction}
\label{sec:supp_pred}

\subsection{The Visibility Audit in Full}

Section~4.2 of the manuscript states that every historical feature is recomputed by
streaming the log in availability order and reading each entity's state as of the query
instant, and that \devnum{0} rows disagree in any of the eight configurations. This is
what that test is.

On \devnum{3{,}000} randomly drawn development jobs all \devnum{34} history and
relational features are recomputed from scratch under each configuration's own
visibility rule and compared against the pipeline's values. The same test applied to the
natural but wrong construction, in which a feature may use any record with a smaller
index, fails on \devnum{97.67\%} of rows under the primary rule, so the test does
discriminate. The graph pipeline is checked separately: an index assertion over
\devnum{863{,}149} rows finds \devnum{0} timeline entries pointing at or after the read,
brute-force recomputation on \devnum{600} rows matches on all \devnum{2{,}400} root
states and all \devnum{2{,}400} neighbour sets, and perturbing or deleting every record
after a cut time leaves all earlier predictions unchanged while a positive control moves
the later ones.

\subsection{Predictors on the Non-Judging Traces}

Table~\ref{tab:s_pred_cross} is the cross-trace counterpart of Table~1 of the
manuscript. The three figures Section~4.4 quotes from it are the serverless STATIC
Spearman correlation and the two CI pools' ENT AUROC; the AUROC column of Table~11 of
the manuscript carries the figure the scheduling results there are read against, which
is the ALL row of each trace.

\begin{table}[H]
\centering
\caption{The same protocol on the non-judging traces, by feature group. ENT uses the
entity's own visible history, STATIC uses only what is declared before the job runs, ALL
is both. RMSE and Spearman are on $\log(1+C)$.}
\label{tab:s_pred_cross}
\small
\begin{tabular}{llcccc}
\hline
Trace & Features & AUROC & AUPRC & RMSE & Spearman \\
\hline
Serverless invocations & ENT    & \devnum{0.9828} & \devnum{0.816} & \devnum{0.3401} & \devnum{0.953} \\
                       & STATIC & \devnum{0.6238} & \devnum{0.084} & \devnum{0.8088} & \devnum{0.021} \\
                       & ALL    & \devnum{0.9847} & \devnum{0.830} & \devnum{0.3185} & \devnum{0.932} \\
CI pool, hardware A    & ENT    & \devnum{0.6334} & \devnum{0.157} & \devnum{0.3964} & \devnum{0.833} \\
                       & STATIC & \devnum{0.4990} & \devnum{0.045} & \devnum{0.7659} & \devnum{0.277} \\
                       & ALL    & \devnum{0.6512} & \devnum{0.177} & \devnum{0.3956} & \devnum{0.837} \\
CI pool, hardware B    & ENT    & \devnum{0.9968} & \devnum{0.984} & \devnum{0.0887} & \devnum{0.985} \\
                       & STATIC & \devnum{0.6915} & \devnum{0.077} & \devnum{0.6411} & \devnum{0.500} \\
                       & ALL    & \devnum{0.9970} & \devnum{0.985} & \devnum{0.0860} & \devnum{0.988} \\
Compute farm           & ENT    & \devnum{0.8903} & \devnum{0.265} & \devnum{1.4214} & \devnum{0.883} \\
                       & STATIC & \devnum{0.8704} & \devnum{0.242} & \devnum{2.2559} & \devnum{0.750} \\
                       & ALL    & \devnum{0.8902} & \devnum{0.260} & \devnum{1.3485} & \devnum{0.892} \\
\hline
\end{tabular}
\end{table}
% Sources of Table~\ref{tab:s_pred_cross} (carried over from 04_prediction.tex, where this
%   table was Table 2 until the sixth revision pass):
%   Serverless invocations and compute farm: evidence/cross_domain/out_cross_domain_table.txt,
%     section B ("predictability of the cost from information at arrival"), rows azure */netbatch *.
%   CI pools A and B: evidence/firefox_ci/out_SUMMARY.txt, section C; A =
%     releng-hardware/gecko-t-osx-1500-m4, B = releng-hardware/gecko-t-linux-talos-2404.
%   The identity variance shares (0.867 serverless, 0.638 compute farm) quoted in Section 4.4
%     are the eta^2 column of the same table in out_cross_domain_table.txt.

%=================================================================
\section{Traces and Workload Characterisation}
\label{sec:supp_data}

\subsection{The Field Audit of the Main Judging Trace}

Section~7.2 of the manuscript states four readings of the main trace. This is the
evidence for them, and the two archive defects that follow from reading the trace this
way.

\paragraph{The timestamp is a result stamp.}
Among later-period runs the gap to the same user's previous event grows one for one with
the cost (slope \devnum{0.99} of the median gap on $C$) while the gap to the next event
barely moves (slope \devnum{0.16}), and the feedback after a user kills a run arrives
within about \devnum{0.1} s whatever its length. A result stamp predicts both readings; a
receipt stamp neither.

\paragraph{The whole-second leak.}
Stamps are recorded to the whole second while costs are recorded to the microsecond, so
the fractional part of a job's cost appears in the gap between its own stamp and the next
event, and a feature like ``seconds since the user's previous record'' reads the label off
the clock. We close it by replacing $t$ with $t + u$, where $u \in [0,1)$ is drawn
deterministically from a fixed key and the record's identifier.

\paragraph{The cost field and the limit.}
The cost field is wall-clock occupancy of one evaluation including interpreter start-up,
read from the 10th percentile of correct submissions (\devnum{0.12}--\devnum{0.18} s
across the two periods) and from its weak dependence on the number of test cases (median
\devnum{$0.196 + 0.007 n$} s in the later period). The platform limit is \devnum{60} s:
the last run to reach it ends at \devnum{2019-08-28} \devnum{10{:}32}, after which the
largest observed cost is \devnum{59.97} s, with \devnum{13} runs in $[59,60)$ s and none
above. Before that date there was no fixed limit, the longest run is \devnum{8{,}901} s,
and the \devnum{1{,}189} runs that end with a kill are almost all above \devnum{360.8} s,
pointing at a container watchdog rather than judging demand.

\paragraph{Co-ending groups, zero-cost blocks and two short archives.}
Groups of runs by the same user that end in the same second are \devnum{7.53\%} of blocks
and \devnum{26.28\%} of capped work; every prediction table in the manuscript is repeated
on rows outside such groups. A further \devnum{1{,}671} blocks record a cost of exactly
zero with empty outputs; they leave the simulated trace and stay in the features, the
labels and every total. The \devnum{2019-2} archive stops at \devnum{2019-09-19}, months
before that term's assessments end, and \devnum{2020-1} is short by about \devnum{3\%} of
its blocks; both serve per-submission modelling and are excluded from anything that treats
a term as a whole.

\paragraph{Against the published release statistics.}
We match the class count in \devnum{18} of \devnum{18} terms and the student and exercise
counts in \devnum{16} of \devnum{18}, and parse \devnum{4.08\%} fewer execution blocks
than the official column sums, whose printed totals disagree with those same columns by
\devnum{$-248$} exercises and \devnum{$+246{,}246$} blocks. We quote the columns.

\subsection{The Continuous-Integration Refit: Sharpness and Circularity}

Section~7.3 of the manuscript reports the held-out refit of the effective capacity
$k_{\mathrm{eff}}$ and the setup term on the two Firefox CI pools. Two checks on that fit
sit here.

The objective is sharp in $k$ rather than flat: on a step-of-one grid only \devnum{5} to
\devnum{8} values of $k$ put the simulated mean wait within \devnum{10\%} of the recorded
one, so the capacity the fit chooses is not one of many equally good answers.

The circularity of the original whole-window fit can be read off the two pools directly.
On pool~A the whole-window and held-out rows are almost identical, because every split
picks nearly the same capacity; on pool~B the reverse split's held-out capacity gives a
mean ratio of \devnum{1.14} on its test part where the whole-window capacity gives
\devnum{0.99}. The difference between those two numbers is the size of the circularity
that fitting on the reported window would have hidden.
% Sources of Section~\ref{sec:supp_data} (carried over from 07_data.tex):
%   Field audit: evidence/codebench_audit/out_tail_and_counts.txt and
%     evidence/codebench_service_v2/out_service_v2.txt (slopes, kill feedback, the 60 s
%     limit and its last over-limit run, the zero-cost blocks, the two short archives and
%     the release-statistics comparison).
%   Sharpness and circularity: evidence/firefox_ci_holdout/out_SUMMARY.txt, section E
%     (5-8 values of k inside 10% of the recorded mean) and section F (the circularity
%     comparison 1.14 vs 0.99 on the talos reverse split); the same rows are in holdout.csv.

%=================================================================
\section{Experiments}
\label{sec:supp_exp}

\subsection{What the Selection Rule Chose}

Table~5 of the manuscript gives the three pre-stated grids, the feasibility rule, the
objective and the tie-break, and names the shape the rule selected at each promise.
Table~\ref{tab:s_setup} gives the selected points in full and the nine per-shape winners.
The same values are in \texttt{protocol\_lock.draft.json}, which is the artefact the
selection was frozen into.

\begin{table}[H]
\centering
\caption{The three selected budget points and the nine per-shape grid winners. Budgets in
work scale as $k/4$; the harm limit is $G/2$, the feasibility threshold of the selection
rule. ``Feasible'' counts the grid points that met it in all \devnum{15} validation cells.}
\label{tab:s_setup}
\small
\begin{tabular}{p{0.20\textwidth}p{0.72\textwidth}}
\toprule
Item & Value \\
\midrule
Selected, $G = \devnum{300}$ s & age-relative: $B_0 = \devnum{60}k/4$ s, $\eta = \devnum{0.50}$, $\gamma = \devnum{0}$; worst-cell gap \devnum{0.4251}, worst-cell harm \devnum{122.8} s against a \devnum{150} s limit, \devnum{44} of \devnum{99} candidate points feasible \\
Selected, $G = \devnum{600}$ s & age-relative: $B_0 = \devnum{120}k/4$ s, $\eta = \devnum{0.75}$, $\gamma = \devnum{0}$; worst-cell gap \devnum{0.7216}, worst-cell harm \devnum{293.3} s against \devnum{300} s, \devnum{70} of \devnum{107} feasible \\
Selected, $G = \devnum{1200}$ s & queue-length: $B_0 = \devnum{0}$, $\eta = \devnum{0}$, $\gamma = \devnum{16}k/4$ s; worst-cell gap \devnum{0.7283}, worst-cell harm \devnum{145.3} s against \devnum{600} s, \devnum{95} of \devnum{111} feasible \\
Per-shape winners & constant: $B_0 = \devnum{329.671}$, \devnum{981.126}, \devnum{2{,}029.958}\,$k/4$ s at the three promises; age-relative: $(\devnum{60}, \devnum{0.50})$, $(\devnum{120}, \devnum{0.75})$, $(\devnum{360}, \devnum{0.75})$; queue-length: $\gamma = \devnum{16}k/4$ s with $B_0 = \devnum{0}$, \devnum{120}$k/4$ s, \devnum{0} \\
\bottomrule
\end{tabular}
\end{table}
% Source of Table~\ref{tab:s_setup}: these are the four rows the manuscript's experimental
%   configuration table (tab:setup, Table 5) carried until the sixth revision pass (referee
%   item G7).  They come from
%   protocol_lock.draft.json (selection_rule, selected_parameters, family_best) and
%   docs/repo_build_status.md third pass, section 3.  Produced by
%   `python scripts/emit_paper_tables.py` into outputs/paper_tables/tab_setup.tex.

\subsection{The Two Budget Shapes the Main Table No Longer Prints}

Section~8.4 of the manuscript compares the three budget shapes at a fixed promise through
the paired differences, which is how the comparison is read. Table~\ref{tab:s_guard_abl}
gives the age-relative and queue-length rows themselves, in the format of Table~8 and in
the same cells; the Guard($G$), Guard-fixed($G$), Fixed($G$) and Skip($G$) rows stay
there. One of the three shape winners coincides with Guard($G$) at each promise, which is
why some rows below repeat a row of Table~8.

\begin{table}[H]
\begin{adjustwidth}{-\extralength}{0cm}
\centering
\caption{The age-relative and queue-length winners at three promises and three loads, in
the format of Table~8 of the manuscript. Excess and harm are the worst over the five
overlays, the rest are means; budgets scale as $k/4$.}
\label{tab:s_guard_abl}
\small
\setlength{\tabcolsep}{3pt}
\begin{tabular}{llrrlrrrr}
\toprule
Load & Policy & Promise & p99$_{\mathrm{dl}}$ & Gap closed & Red.\ (\%) & Max exc. & Harm & Fired (\%) \\
\midrule
\multicolumn{9}{l}{$\rho = 0.5$, $k = 7.8$} \\
 & Guard-age($\devnum{300}$)  & \devnum{300} & \devnum{15.58} & \devnum{0.705 [0.678, 0.724]} & \devnum{43.7} & \devnum{192.5} & \devnum{132.1} & \devnum{1.35} \\
 & Guard-queue($\devnum{300}$) & \devnum{300} & \devnum{15.27} & \devnum{0.723 [0.697, 0.742]} & \devnum{44.8} & \devnum{193.4} & \devnum{94.3} & \devnum{21.60} \\
 & Guard-age($\devnum{600}$)  & \devnum{600} & \devnum{14.97} & \devnum{0.740 [0.713, 0.760]} & \devnum{45.8} & \devnum{492.2} & \devnum{237.0} & \devnum{0.17} \\
 & Guard-queue($\devnum{600}$) & \devnum{600} & \devnum{14.88} & \devnum{0.744 [0.718, 0.765]} & \devnum{46.1} & \devnum{504.5} & \devnum{106.0} & \devnum{0.16} \\
 & Guard-age($\devnum{1200}$) & \devnum{1200} & \devnum{14.73} & \devnum{0.751 [0.723, 0.773]} & \devnum{46.5} & \devnum{1{,}049.1} & \devnum{237.0} & \devnum{0.01} \\
 & Guard-queue($\devnum{1200}$) & \devnum{1200} & \devnum{14.84} & \devnum{0.746 [0.719, 0.768]} & \devnum{46.2} & \devnum{907.8} & \devnum{94.3} & \devnum{22.01} \\
\midrule
\multicolumn{9}{l}{$\rho = 0.8$, $k = 5$} \\
 & Guard-age($\devnum{300}$)  & \devnum{300} & \devnum{55.94} & \devnum{0.738 [0.697, 0.761]} & \devnum{53.0} & \devnum{209.7} & \devnum{114.1} & \devnum{11.57} \\
 & Guard-queue($\devnum{300}$) & \devnum{300} & \devnum{54.69} & \devnum{0.753 [0.713, 0.777]} & \devnum{54.0} & \devnum{210.7} & \devnum{117.6} & \devnum{17.53} \\
 & Guard-age($\devnum{600}$)  & \devnum{600} & \devnum{47.47} & \devnum{0.837 [0.815, 0.851]} & \devnum{60.1} & \devnum{507.8} & \devnum{319.6} & \devnum{1.31} \\
 & Guard-queue($\devnum{600}$) & \devnum{600} & \devnum{47.69} & \devnum{0.834 [0.814, 0.847]} & \devnum{59.9} & \devnum{508.5} & \devnum{117.7} & \devnum{1.35} \\
 & Guard-age($\devnum{1200}$) & \devnum{1200} & \devnum{45.16} & \devnum{0.864 [0.844, 0.876]} & \devnum{62.0} & \devnum{1{,}102.9} & \devnum{370.5} & \devnum{0.29} \\
 & Guard-queue($\devnum{1200}$) & \devnum{1200} & \devnum{46.85} & \devnum{0.844 [0.823, 0.857]} & \devnum{60.6} & \devnum{1{,}104.9} & \devnum{117.6} & \devnum{9.93} \\
\midrule
\multicolumn{9}{l}{$\rho = 1.0$, $k = 4$} \\
 & Guard-age($\devnum{300}$)  & \devnum{300} & \devnum{163.40} & \devnum{0.478 [0.387, 0.602]} & \devnum{40.2} & \devnum{229.1} & \devnum{134.8} & \devnum{28.64} \\
 & Guard-queue($\devnum{300}$) & \devnum{300} & \devnum{169.22} & \devnum{0.453 [0.368, 0.607]} & \devnum{38.1} & \devnum{231.4} & \devnum{144.8} & \devnum{30.99} \\
 & Guard-age($\devnum{600}$)  & \devnum{600} & \devnum{89.24} & \devnum{0.801 [0.764, 0.842]} & \devnum{67.3} & \devnum{523.7} & \devnum{339.9} & \devnum{7.15} \\
 & Guard-queue($\devnum{600}$) & \devnum{600} & \devnum{91.09} & \devnum{0.793 [0.758, 0.833]} & \devnum{66.6} & \devnum{515.8} & \devnum{123.4} & \devnum{7.08} \\
 & Guard-age($\devnum{1200}$) & \devnum{1200} & \devnum{72.52} & \devnum{0.874 [0.851, 0.893]} & \devnum{73.4} & \devnum{1{,}117.3} & \devnum{453.3} & \devnum{0.81} \\
 & Guard-queue($\devnum{1200}$) & \devnum{1200} & \devnum{82.67} & \devnum{0.830 [0.801, 0.857]} & \devnum{69.7} & \devnum{1{,}110.4} & \devnum{144.8} & \devnum{6.21} \\
\bottomrule
\end{tabular}
\end{adjustwidth}
\end{table}
% Source of Table~\ref{tab:s_guard_abl}: the eighteen rows the manuscript's guard table
%   (tab:guard, Table 8) carried until the sixth revision pass (referee item G1).
%   outputs/paper_tables/tab_guard.tex,
%   produced by `python scripts/emit_paper_tables.py` (see GENERATED.md), from
%   outputs/dev_tables/main_table.csv, rows Guard-age / Guard-queue at each promise and
%   levels 0, 1, 2; columns p99_dl_s, gap_closed with gap_closed_lo/hi, reduction_pct,
%   max_excess_s, harm_s, fired_pct.  Nothing was recomputed for this file.

\subsection{The Residual on All Nine Cells}

Table~10 of the manuscript prints the three Guard($\devnum{600}$) cells and one reversed
score. Table~\ref{tab:s_resid} gives all nine. The reading Section~8.5 takes from them is
that the residual is a property of the schedule and barely moves with the policy: inside
each load the ratios are similar in scale, but are not equal to the third decimal.
An $R^2$ rounded to unity is likewise not exact equality.

\begin{table}[H]
\centering
\caption{The residual $D_i = k(W_P[i] - W_{\mathrm{FCFS}}[i]) - (\mathrm{In}_i -
\mathrm{Out}_i)$ of Theorem~1 of the manuscript on overlay~0,
\devnum{17{,}634{,}760} jobs per cell. Ratio is $\max|D_i|$ over the bound $2(k-1)L$,
which is \devnum{14}$L$, \devnum{8}$L$ and \devnum{6}$L$ at the three loads;
$R^2$ and the worst error are for $(\mathrm{In}_i - \mathrm{Out}_i)/k$
against the realised excess over first-come first-served.}
\label{tab:s_resid}
\small
\setlength{\tabcolsep}{3pt}
\begin{tabular}{llrrrrr}
\toprule
Load & Policy & $\max|D_i|/L$ & Ratio & $D_i = 0$ & $R^2$ & worst err.\ (s) \\
\midrule
$\rho = \devnum{0.5}$, $k = \devnum{8}$
 & SPJF-E              & \devnum{6.788} & \devnum{0.485} & \devnum{95.6\%} & \devnum{0.985} & \devnum{50.9} \\
 & Guard($\devnum{600}$) & \devnum{6.788} & \devnum{0.485} & \devnum{95.6\%} & \devnum{0.983} & \devnum{50.9} \\
 & Reversed score      & \devnum{6.788} & \devnum{0.485} & \devnum{95.6\%} & \devnum{0.996} & \devnum{50.9} \\
\midrule
$\rho = \devnum{0.8}$, $k = \devnum{5}$
 & SPJF-E              & \devnum{3.922} & \devnum{0.490} & \devnum{87.2\%} & \devnum{0.996} & \devnum{47.1} \\
 & Guard($\devnum{600}$) & \devnum{3.922} & \devnum{0.490} & \devnum{87.2\%} & \devnum{0.995} & \devnum{47.1} \\
 & Reversed score      & \devnum{3.881} & \devnum{0.485} & \devnum{87.5\%} & \devnum{0.999} & \devnum{46.6} \\
\midrule
$\rho = \devnum{1.0}$, $k = \devnum{4}$
 & SPJF-E              & \devnum{3.082} & \devnum{0.514} & \devnum{81.1\%} & \devnum{0.998} & \devnum{46.2} \\
 & Guard($\devnum{600}$) & \devnum{3.082} & \devnum{0.514} & \devnum{81.1\%} & \devnum{0.998} & \devnum{46.2} \\
 & Reversed score      & \devnum{2.994} & \devnum{0.499} & \devnum{81.7\%} & \devnum{1.000} & \devnum{44.9} \\
\bottomrule
\end{tabular}
\end{table}
% Source of Table~\ref{tab:s_resid}: outputs/paper_tables/tab_resid.tex, produced by
%   `python scripts/emit_paper_tables.py` (see GENERATED.md), from
%   outputs/dev_tables/identity_residuals.csv, rows overlay = 0, levels 0/1/2, policies
%   SPJF-E / Guard(600) / SPJF-reversed; columns max_abs_residual_L, ratio_to_bound,
%   share_residual_zero, r2_net_over_k, max_abs_error_s.  The load labels are the paper's
%   (rho = 0.5 / 0.8 / 1.0); the package prints the level index there.  This is the table
%   the manuscript's residual table (tab:resid, Table 10) printed in full until the sixth
%   revision pass (referee item G10); it now keeps four of these nine rows.

\subsection{A Single Server}

Section~8.6 of the manuscript reports the single-server configuration through five gap
figures, three harm ratios and two realised excesses. Table~\ref{tab:s_k1} is the
configuration in full. The parameters are the ones the selection rule chose on the
\devnum{15} multi-server validation cells; $k = 1$ is not among them, and we do not
re-tune on it.

The firing rate of the three unguarded policies is printed \verb|---|: the column counts
dispatch epochs at which the guard overrides the base policy, and these three runs have no
guard. Scoring first-come first-served at \devnum{100\%} by the convention that every
first-come first-served dispatch serves the queue head, as an earlier version of this
table did, compares it with the guarded runs on a quantity it does not have.

\begin{table}[H]
\centering
\caption{Single server, busy-hour $\rho = \devnum{0.80}$, \devnum{3{,}209{,}347} jobs. Waits and
excesses in seconds. The parameters are the ones the rule selected on the multi-server
validation cells; this configuration is not among them. The skip count is
$\kappa = \lfloor Gk/L \rfloor - (2k-2)$.}
\label{tab:s_k1}
\small
\setlength{\tabcolsep}{3pt}
\begin{tabular}{lrrlrrr}
\toprule
Policy & Promise & p99$_{\mathrm{dl}}$ & Gap closed & Max exc. & Harm & Fired (\%) \\
\midrule
FCFS                         & ---           & \devnum{221.27} & \devnum{0.000} & \devnum{0.0} & \devnum{0.0} & --- \\
SJF on true sizes            & ---           & \devnum{45.32} & \devnum{1.000} & \devnum{2{,}054.5} & \devnum{777.9} & --- \\
SPJF-E                       & ---           & \devnum{60.91} & \devnum{0.911 [0.854, 0.933]} & \devnum{2{,}520.3} & \devnum{656.4} & --- \\
Guard($\devnum{300}$)        & \devnum{300} & \devnum{89.96} & \devnum{0.746 [0.571, 0.829]} & \devnum{297.9} & \devnum{81.8} & \devnum{8.17} \\
Guard-fixed($\devnum{300}$)  & \devnum{300} & \devnum{125.68} & \devnum{0.543 [0.207, 0.718]} & \devnum{140.3} & \devnum{119.5} & \devnum{17.09} \\
Guard-age($\devnum{300}$)    & \devnum{300} & \devnum{89.96} & \devnum{0.746 [0.571, 0.829]} & \devnum{297.9} & \devnum{81.8} & \devnum{8.17} \\
Guard-queue($\devnum{300}$)  & \devnum{300} & \devnum{105.15} & \devnum{0.660 [0.532, 0.715]} & \devnum{297.9} & \devnum{53.9} & \devnum{29.15} \\
Fixed($\devnum{300}$)        & \devnum{300} & \devnum{83.18} & \devnum{0.785 [0.584, 0.865]} & \devnum{297.9} & \devnum{246.3} & \devnum{6.33} \\
Skip($\devnum{300}$), $\kappa=5$  & \devnum{300} & \devnum{220.39} & \devnum{0.005 [-0.002, 0.019]} & \devnum{95.1} & \devnum{81.8} & \devnum{73.17} \\
Guard($\devnum{600}$)        & \devnum{600} & \devnum{71.07} & \devnum{0.854 [0.780, 0.895]} & \devnum{595.5} & \devnum{165.5} & \devnum{2.53} \\
Guard-fixed($\devnum{600}$)  & \devnum{600} & \devnum{82.10} & \devnum{0.791 [0.602, 0.868]} & \devnum{303.0} & \devnum{246.3} & \devnum{6.16} \\
Guard-age($\devnum{600}$)    & \devnum{600} & \devnum{71.07} & \devnum{0.854 [0.780, 0.895]} & \devnum{595.5} & \devnum{165.5} & \devnum{2.53} \\
Guard-queue($\devnum{600}$)  & \devnum{600} & \devnum{91.52} & \devnum{0.737 [0.640, 0.812]} & \devnum{582.0} & \devnum{81.8} & \devnum{3.43} \\
Fixed($\devnum{600}$)        & \devnum{600} & \devnum{68.48} & \devnum{0.868 [0.783, 0.906]} & \devnum{593.4} & \devnum{563.6} & \devnum{2.29} \\
Skip($\devnum{600}$), $\kappa=10$ & \devnum{600} & \devnum{219.89} & \devnum{0.008 [-0.004, 0.037]} & \devnum{100.9} & \devnum{83.2} & \devnum{66.98} \\
Guard($\devnum{1200}$)       & \devnum{1200} & \devnum{98.00} & \devnum{0.701 [0.640, 0.742]} & \devnum{1{,}157.4} & \devnum{53.9} & \devnum{26.15} \\
Guard-fixed($\devnum{1200}$) & \devnum{1200} & \devnum{69.03} & \devnum{0.865 [0.789, 0.905]} & \devnum{564.2} & \devnum{509.3} & \devnum{2.57} \\
Guard-age($\devnum{1200}$)   & \devnum{1200} & \devnum{63.80} & \devnum{0.895 [0.835, 0.920]} & \devnum{1{,}182.3} & \devnum{366.4} & \devnum{0.68} \\
Guard-queue($\devnum{1200}$) & \devnum{1200} & \devnum{98.00} & \devnum{0.701 [0.640, 0.742]} & \devnum{1{,}157.4} & \devnum{53.9} & \devnum{26.15} \\
Fixed($\devnum{1200}$)       & \devnum{1200} & \devnum{63.19} & \devnum{0.898 [0.837, 0.923]} & \devnum{1{,}180.4} & \devnum{656.4} & \devnum{0.45} \\
Skip($\devnum{1200}$), $\kappa=20$ & \devnum{1200} & \devnum{217.64} & \devnum{0.021 [-0.008, 0.084]} & \devnum{168.0} & \devnum{97.7} & \devnum{58.67} \\
\bottomrule
\end{tabular}
\end{table}
% Source of Table~\ref{tab:s_k1}: outputs/paper_tables/tab_k1.tex, produced by
%   `python scripts/emit_paper_tables.py` (see GENERATED.md), from
%   outputs/dev_tables/k1/main_table.csv; columns p99_dl_s, gap_closed with gap_closed_lo/hi,
%   max_excess_s, harm_s, fired_pct, and the promise G.  The harm ratios 3.01 / 3.41 / 12.18 in
%   Section 8.6 are Fixed(G) harm over Guard(G) harm from the same rows.  The adversarial excess
%   597.0 s is the Guard(600)-top1short row of the same file, column max_excess_s.  The two
%   third-decimal double-rounding errors the verifier found in an earlier version of this
%   table (item G2) are gone with the rows and the rounding: every figure here is printed once
%   from the package's own full-precision column.  This was the manuscript's single-server
%   table until the sixth revision pass (referee item G2); the three unguarded firing rates that the
%   package prints as 100.00 / 0.00 / 0.00 are printed --- here (referee item RR-24), which is
%   the one cell where this table differs from the package's own.

%=================================================================
\section{Deferred Proofs}
\label{sec:supp_proofs}

Appendices~A.1 and~A.5 of the manuscript give the ideas behind the arguments
proved here. The workload comparison uses continuity and left limits, whereas
consistency is proved by induction. The tightness constructions remain in
Appendices~A.2--A.4, A.6 and~A.7 of the manuscript. We use the manuscript's notation:
$U_Q(t)$ is the unfinished work at time $t$ under policy $Q$, and $N_Q(t)$ is the
number of jobs present.

\subsection{The Workload Comparison}
\label{sec:supp_gap}

\begin{proof}[Proof of Lemma~1 of the manuscript]
Put $D = U_{Q_1} - U_{Q_2}$. Arrivals add the same jump to both terms, so $D$ is
continuous, and between arrivals it changes at rate
$-(\min(N_{Q_1},k) - \min(N_{Q_2},k))$; it is therefore piecewise linear with
breakpoints among the finitely many arrival and completion instants of the two
runs.

If $U_{Q_1}(u) > (k-1)L$ then more than $k-1$ jobs are present under $Q_1$, each
carrying at most $L$, so every server of $Q_1$ is busy and
$\min(N_{Q_1},k) = k \ge \min(N_{Q_2},k)$: $D$ does not increase at $u$. Fix $t$
and let
$S = \{u \le t : U_{Q_1}(u) \le (k-1)L\}$, which is non-empty because $U_{Q_1}$
vanishes
before the first arrival, and put $t_0 = \sup S$. The supremum need not be
attained, since $U_{Q_1}$ jumps upwards at arrivals, so we use the left limit:
there are $u_n \uparrow t_0$ with $U_{Q_1}(u_n) \le (k-1)L$, hence
$U_{Q_1}(t_0^-) \le (k-1)L$ and, $D$ being continuous,
$D(t_0^-) \le U_{Q_1}(t_0^-) \le (k-1)L$ because $U_{Q_2} \ge 0$. On $(t_0, t]$
we have
$U_{Q_1} > (k-1)L$, so $D$ does not increase there, and $D(t) \le (k-1)L$.
Exchanging $Q_1$ and $Q_2$ gives the other direction. That the constant cannot be
lowered is the cascade of Appendix~A.2 of the manuscript, which reaches
$(k-1)L\,(1 - ((k-1)/k)^m)$ after $m$ rounds. At $k = 1$ the bound reads
$|D| \le 0$ and is attained everywhere: one server executes work at rate $1$
whenever any job is present, so $U_{Q_1} \equiv U_{Q_2}$.
\end{proof}

\begin{proof}[Proof of the strict inequality in Lemma~1, for $k \ge 2$]
Keep the notation above, so $D = U_{Q_1} - U_{Q_2}$ is continuous and piecewise
linear with finitely many pieces. Suppose $D(t) = (k-1)L$ for some $t$ and let
$t_1$ be the smallest such $t$; it exists because $\{D = (k-1)L\}$ is closed and
$D(0) = 0 < (k-1)L$, which uses $L > 0$ and $k \ge 2$.

Let $I$ be the linear piece immediately to the left of $t_1$, which has positive
length. Its slope is strictly positive: a negative slope would put
$D > (k-1)L$ somewhere on $I$ and a zero slope would put $D = (k-1)L$ on all of
$I$, and either contradicts the minimality of $t_1$. The slope of $D$ is
$-(\min(N_{Q_1},k) - \min(N_{Q_2},k))$, so $\min(N_{Q_1},k) < \min(N_{Q_2},k)$
on $I$ and
therefore $N_{Q_1} \le k-1$ there.

With at most $k-1$ jobs present and $Q_1$ work-conserving, none of $Q_1$'s jobs
waits, so all of them are in service; hence $U_{Q_1}$ falls at rate $N_{Q_1}$ on
$I$,
and $U_{Q_1} \le (k-1)L$ on $I$ because each of the at most $k-1$ present jobs
carries at most $L$. Taking left limits at $t_1$ and using the continuity of
$D$,
\[
  U_{Q_1}(t_1^-) - U_{Q_2}(t_1^-) = (k-1)L , \qquad U_{Q_1}(t_1^-) \le (k-1)L ,
  \qquad U_{Q_2} \ge 0 ,
\]
which forces $U_{Q_2}(t_1^-) = 0$ and $U_{Q_1}(t_1^-) = (k-1)L$. Since $L > 0$
and $k \ge 2$, this
gives $N_{Q_1} \ge 1$ on $I$ --- with $N_{Q_1} = 0$ we would have
$U_{Q_1}(t_1^-) = 0 = (k-1)L$ ---
so $U_{Q_1}$ is strictly decreasing on $I$ and
$U_{Q_1}(t_1^-) < \sup_I U_{Q_1} \le (k-1)L$, a contradiction. Exchanging $Q_1$
and $Q_2$ gives the other direction. Both steps that invoke $(k-1)L > 0$ fail at
$k = 1$, where the statement is false: there $U_{Q_1} \equiv U_{Q_2}$ and the
bound $0$ is attained at every instant.
\end{proof}

\subsection{Consistency, and Why the Converse Fails}
\label{sec:supp_consistency}

\begin{proof}[Proof of Proposition~1 of the manuscript]
Induct on dispatch epochs: up to the first epoch at which the two runs could
differ their completion histories agree, hence so does $\mathrm{over}[\cdot]$,
so $E(t)$ is empty in the wrapped run and it copies $A$. The first consequence
uses $\mathrm{over}[q](t) \le \mathrm{In}^A_q$, a completed overtaker of $q$
having been dispatched while $q$ waited; the second uses Corollary~2 of the
manuscript.
\end{proof}

The converse fails, so this is not an equivalence. Take $k = 1$, two jobs of
service $1$ at $t = 0$, the base policy ``higher rank first'', and $B = 1$. Both
runs dispatch in the order $(1,0)$, yet at the second dispatch the waiting job
has $\mathrm{over} = 1 \ge B$, so the guard fires and happens to choose what the
base would have chosen. Nor is there multiplicative consistency: with $k = 1$,
one job of service $L$ followed by $m$ jobs of service $1$ all at $t = 0$, the
wrapper with $B = 0$ is first-come first-served while the base ordering by
service time is not, and the ratio of mean waits is $(2L + m - 1)/(m+1)$,
unbounded in $L$.
% Source of Section~\ref{sec:supp_proofs}: sections/A_proofs.tex, Appendices A.1 and A.5
%   of the manuscript, moved here in the seventh revision pass under referee item G3.
%   The two arguments are verbatim; only the cross-references were rewritten, since a
%   \ref into the manuscript does not resolve from this file: Lemma~\ref{lem:gap} is
%   Lemma 1, Proposition~\ref{prop:consistency} is Proposition 1,
%   Corollary~\ref{cor:necessary} is Corollary 2 and Appendix~\ref{app:cascade} is
%   Appendix A.2.  Simulated and checked in evidence/guard_theory; nothing was
%   recomputed for this file.

\section{Reproducibility and Provenance of the Empirical Results}
\label{sec:supp_repro}

\subsection{Feature Definitions, Targets and Rolling Splits}

The executable feature dictionary is \path{src/spjf_guard/features/sweep.py}:
\path{HIST_COLS}, \path{CTX_COLS}, \path{GROUPS} and
\path{feature_frame} specify the entity aggregates, missing values and event
order; \path{src/spjf_guard/data/events.py} specifies code columns and clock inputs.
M4 combines static code, assessment/calendar context, exercise history, user history,
user--exercise history and previous-event history. It excludes the relational
extensions. Outcome aggregates contain log-cost moments and quantiles, error/heavy
rates and recent-result summaries; counts and elapsed time since arrival are known
at arrival and are distinguished from outcome fields. The rolling-history window,
tie order and update formulas are defined in that dictionary. Static uses only the
code and context blocks. The independent source-clock audit in Section~S1 verifies
that dictionary under its original clock, not under an arbitrary replayed policy.

Package scores are fitted by \path{src/spjf_guard/predict/forward.py}. For each
target term, training includes only terms whose first arrival precedes that target,
and only records whose result is readable strictly before its first arrival.
History thresholds and user-slowness cutpoints are refitted from those training
rows; a model is then frozen for the target term. The targets are
\devnum{2020-ERE}, \devnum{2020-2}, \devnum{2021-1}, \devnum{2021-2},
\devnum{2022-1} and \devnum{2022-2}. Starting at \devnum{2018-1}, their training
periods end respectively at \devnum{2019-2}, \devnum{2020-1}, \devnum{2020-2},
\devnum{2021-1}, \devnum{2021-2} and \devnum{2022-1}; availability filtering still
applies at each cutoff. Validation overlays exclude the final development term
from selection, but that term and the earlier family exploration are development
data, not sealed testing.
% Sources: outputs/dev_predictor/manifest.json, arguments.targets;
% evidence/ranking_score/fit_targets.csv, target and train_sems; outputs/selection_v3/
% manifest.json, validation pool; prechecks/heldout_scope/heldout_scope.md, Q2.

The score objectives are in \path{src/spjf_guard/predict/scores.py}, and the
deterministic fit settings in \texttt{lgb\_parameters} and the configuration
identified by the output manifest. \texttt{spjf\_e} is a Tweedie regression of
capped seconds; \texttt{spjf\_log} is squared-error regression of log-transformed
cost. In contrast, the earlier M4 AUROC \devnum{0.9283 [0.9118, 0.9415]} is the
heavy-class classifier probability, averaged over \devnum{3} seeds, trained on
\devnum{2018-1}--\devnum{2019-2} and tested only on \devnum{2022-2}. Those
retrospective classifier predictions never enter scheduling. The package's pooled
\texttt{spjf\_e} AUROC \devnum{0.9360} combines \devnum{402{,}358} unreplicated
forward-prediction rows. Its evaluation heavy threshold stays at the core training
p95, \devnum{1.559043} s; it is not the threshold refitted inside each model's
history features. Predictor intervals resample users, whereas scheduling intervals
resample weeks. The two AUROCs compare different fits, targets and pooling and
must not be read as a measured improvement on a common split.
% Sources: evidence/codebench_service_v2/out_service_v2.txt, SUMMARY 1;
% outputs/dev_predictor/predictor_metrics.csv, pooled/spjf_e; manifest.json,
% heavy_reference_semesters, heavy_threshold_s and bootstrap_block.

\subsection{Exact Selection Grids and Counts}

The full expanded grid, its definitions, counts and deduplication totals are in
\path{outputs/selection_v3/selection_protocol.json}; per-cell values, worst-cell
values and selections are in the adjacent \path{selection_grid.csv},
\path{selection_worst.csv}, \path{selection_frontier.csv} and
\path{selected_parameters.csv}. All budget and queue-work bases below scale
by $k/\devnum{4}$ before dispatch arithmetic. The cap is determined by the promise.
% Source: outputs/selection_v3/selection_protocol.json and
% outputs/dev_tables/policy_parameters.csv, k and B0/gamma values.

The constant grid has the following \devnum{39} base budgets (seconds), shared
across promises:
\begin{quote}\small
\devnum{15}, \devnum{17.99}, \devnum{21.576}, \devnum{25.877}, \devnum{30}, \devnum{31.035}, \devnum{37.222}, \devnum{44.641}, \devnum{53.54}, \devnum{60}, \devnum{64.212}, \devnum{77.012}, \devnum{92.363}, \devnum{110.774}, \devnum{120}, \devnum{132.855}, \devnum{159.338}, \devnum{191.099}, \devnum{229.192}, \devnum{240}, \devnum{274.878}, \devnum{329.671}, \devnum{360}, \devnum{395.386}, \devnum{474.201}, \devnum{480}, \devnum{568.726}, \devnum{600}, \devnum{682.093}, \devnum{818.058}, \devnum{981.126}, \devnum{1176.7}, \devnum{1411.26}, \devnum{1692.57}, \devnum{2029.96}, \devnum{2434.6}, \devnum{2919.9}, \devnum{3501.94}, \devnum{4200}.
\end{quote}
% Source: outputs/selection_v3/selection_protocol.json, expanded_points, fixed with promise_s=-1.

It is the union of \devnum{32} logarithmically spaced values between
\devnum{15} and \devnum{4200} s, rounded as recorded, and the anchors
$\devnum{15},\devnum{30},\devnum{60},\devnum{120},\devnum{240},\devnum{360},
\devnum{480},\devnum{600}$, with duplicate values removed. A full-cap constant
point is added at each of $G=\devnum{300},\devnum{600},\devnum{1200}$ s.
This makes \devnum{42} expanded constant entries.
% Source: outputs/selection_v3/selection_protocol.json, definition.fixed and expanded_points.

For each promise the age-relative grid is the Cartesian product
\[
 B_{0,\mathrm{base}}\in\{\devnum{0},\devnum{15},\devnum{30},\devnum{60},
 \devnum{120},\devnum{240},\devnum{360},\devnum{480},\devnum{600}\},\qquad
 \eta\in\{\devnum{0},\devnum{0.25},\devnum{0.5},\devnum{0.75},
 \devnum{0.9},\devnum{0.95}\},\quad\gamma=0.
\]
The queue-length grid is
\[
 B_{0,\mathrm{base}}\in\{\devnum{0},\devnum{30},\devnum{120}\},\qquad
 \eta\in\{\devnum{0},\devnum{0.9}\},\qquad
 \gamma_{\mathrm{base}}\in\{\devnum{1},\devnum{4},\devnum{16}\}.
\]
Thus the family counts over all promises are
$\devnum{42}+\devnum{162}+\devnum{54}=\devnum{258}$ expanded points.
Deduplicating identical dispatch rules after scaling and saturation leaves
\devnum{243} distinct schedules at each evaluated server count. This is a
computational deduplication, not a reduction of the declared search space.
% Source: outputs/selection_v3/selection_protocol.json, definition.capped/hybrid,
% expanded_points and per_server_count.

Candidate counts per promise answer a different question. The constant entries
shared across promises qualify for a given $G$ only if their worst-cell certified
promise is at most $G$; an age-relative or queue-length entry belongs to its own
declared promise. Feasibility then additionally requires harm $\le G/2$ in every
validation cell. Table~\ref{tab:s_grid_counts} separates these filters, explaining
why the per-promise totals differ from the expanded search size. The objective is
the minimum deadline gap over all validation cells, with ties resolved by smaller
worst harm, then smaller $B_0$, $\eta$ and $\gamma$.

\begin{table}[htbp]
\centering\small
\caption{Candidates / feasible points after promise assignment and the harm filter.}
\label{tab:s_grid_counts}
\begin{tabular}{rrrrr}
\toprule
$G$ (s) & Constant & Age-relative & Queue-length & Joint \\
\midrule
\devnum{300} & \devnum{27} / \devnum{22} & \devnum{54} / \devnum{15} & \devnum{18} / \devnum{7} & \devnum{99} / \devnum{44} \\
% Source: outputs/selection_v3/selection_protocol.json, counts, promise_s=300.
\devnum{600} & \devnum{35} / \devnum{31} & \devnum{54} / \devnum{30} & \devnum{18} / \devnum{9} & \devnum{107} / \devnum{70} \\
% Source: outputs/selection_v3/selection_protocol.json, counts, promise_s=600.
\devnum{1200} & \devnum{39} / \devnum{35} & \devnum{54} / \devnum{43} & \devnum{18} / \devnum{17} & \devnum{111} / \devnum{95} \\
% Source: outputs/selection_v3/selection_protocol.json, counts, promise_s=1200.
\bottomrule
\end{tabular}
\end{table}


The aging comparator has the credit grid
\devnum{0.00001}, \devnum{0.00003}, \devnum{0.0001}, \devnum{0.0003},
\devnum{0.001}, \devnum{0.003}, \devnum{0.01}, \devnum{0.03}, \devnum{0.1},
\devnum{0.3} and \devnum{1}. It uses the same validation pool, harm threshold and
worst-cell objective, with ties to less harm then larger credit. It uses
conservative scores; the guard grids used the original package score. The
selected guard parameters are subsequently held fixed across visibility controls,
so this is not a claim of reselection under conservative or exact histories.
The aging label carries no proved guarantee.
% Sources: outputs/selection_aging/aging_protocol.json and manifest.json;
% outputs/selection_v3/manifest.json, score_array=tweedie, score_parquet=null;
% outputs/dev_visibility/manifest.json, selected_parameters_unchanged.

\subsection{Replication, History and Aggregation Order}

Each overlay is built from the same \devnum{40} class-terms replicated
\devnum{44} times with independent whole-week shifts; each of the five overlays
therefore reuses the same jobs. Original scores are computed on the source calendar
and copied with their source jobs. Conservative histories are namespaced within a
class-term, use only replayed submissions and share the copy's shift. These copies
are not independently trained predictors or independent workload samples.
% Source: evidence/main_v3/out_main_v31.txt, trace construction table;
% prechecks/weakness1_attack/REPORT.md, development capacity experiment.

For the pending exact refinement, a copy means the same class-term, replication
round and shift. Its own submission outcomes follow policy-specific completions.
Other classes and terms are exogenous history in that copy's original relative
clock: when the target copy shifts by $s_i$, an external result is interpreted at
$\mathrm{done}_j+s_i$, not at another copy's independent shift. This preserves
earlier completed history without claiming a single platform whose copied users
share dynamically changing state. A separate sensitivity removes entire records
from other replayed classes in the same term, including arrival counts and elapsed
time since arrival; it then repeats same-copy refinement. Ordinary outcome
withholding preserves already-observed arrivals. The event order releases positive
duration completions before reads at the same instant, and zero-duration outcomes
after those reads. These conventions are specified in
\path{docs/policy_consistent_visibility.md} and implemented by
\path{src/spjf_guard/features/refine.py}.

In the package main and visibility tables, each overlay's deadline p99 is first
computed over its eligible jobs; its gap is then computed against that overlay's
FCFS and SJF. Reported p99 and gap are separate arithmetic means over overlays.
Consequently the displayed mean waits need not algebraically recover the displayed
mean gap. Harm and maximum excess are maxima over overlays. Firing is
queue-weighted within a cell, then averaged across overlays. Each bootstrap draw
uses the same week multiplicities for every policy and overlay, recomputes weighted
job quantiles and cell-level ratios, and averages the ratios. No table uses an
average of weekly p99 values. The capacity-sweep gap instead uses the ratio of
averaged overlay quantiles, as explicitly identified in Section~S6.
% Sources: outputs/dev_tables/main_cells.csv and main_table.csv, per-cell/aggregate
% columns; outputs/dev_visibility/visibility_cells.csv and visibility_comparison.csv;
% prechecks/weakness1_attack/capacity_curve.csv. Algorithms: scripts/run_main.py,
% spjf_guard.experiment.metrics and prechecks/weakness1_attack/summarize_capacity.py.

\subsection{Which Results the Package Reproduces}

The package produces the main CodeBench ranking, guard, adversarial and residual
tables from \path{outputs/dev_tables/}, including its separate single-server
directory; predictor metrics from \path{outputs/dev_predictor/}; and visibility
comparisons from \path{outputs/dev_visibility/}. Guard selection and aging selection
are in \path{outputs/selection_v3/} and \path{outputs/selection_aging/}. Their
manifests identify inputs, settings and output hashes. The preserved original
recipe is \path{configs/main_original_84932d9.yaml}; the original outputs are not
claimed to have been regenerated under an unfinished visibility revision.

\begin{sloppypar}
Earlier predictor-family comparisons and source-clock feature sensitivities are
archived under \path{evidence/codebench_service_v2/} and
\path{evidence/predictor_neural/}. The refit-control and GRU rows in the main
ranking table come from \path{evidence/main_v3/}; the ranking-loss table uses its
own fits in \path{evidence/ranking_score/}. Cross-domain rows combine the package
CodeBench row with \path{evidence/cross_domain/}, \path{evidence/accoding_v2/},
\path{evidence/firefox_ci/} and the LPC-EGEE study below. CI chronological-refit
validation is in \path{evidence/firefox_ci_holdout/}. The capacity and physical
experiments are separate analyses in \path{prechecks/weakness1_attack/}.
\end{sloppypar}

The trace-construction sensitivity table (Table~12 in the unrevised numbering)
uses the earlier source-clock M4 log-scale score, overlays
\devnum{0}--\devnum{2}, and a fixed completed-work budget
$B=\devnum{600}$ s. This is neither a \devnum{600}-second promise nor the selected
age-relative Guard(\devnum{600}). It supports the earlier ordering result's
sensitivity to term inclusion, the alternate timestamp interpretation, and
merging or dropping co-ending runs. It does not test the final expected-cost
guard, conservative history or policy-specific refinement under those altered
inputs. Numerical agreement of the exact-arithmetic package and its independent
reference is also distinct from the archived earlier-pipeline validation.
% Source: evidence/codebench_service_v2/out_service_v2.txt, SUMMARY 5,
% primary(0-2), SPJF-M4 and GUARD-M4-B600 rows; evidence/main_v3_verify/out_VERDICT.txt.

The sealed procedure is specified by \path{docs/sealed_run_procedure.md}. It
produces separate multi-server, single-server, predictor and visibility outputs,
including the policy-specific exact stage, with unchanged declared selection
rules and parameters. It does not extend the earlier predictor-family or
cross-domain studies. No sealed result is present in this supplement.

\section{Capacity, Physical Queueing and Feedback}
\label{sec:supp_physical}

These development experiments address distinct questions. The capacity sweep changes
the size of the constructed pool while preserving its complete offered stream. The
physical experiment checks an executing timed-work service on a selected window.
Neither identifies a historical queue on CodeBench or a causal deployment effect.
The scores are the archived embedded Tweedie predictions, not the later package
refit, conservative score or pending exact visibility variant.

\subsection{Capacity Sensitivity on Complete Traces}

The declared common grid uses every integer $k=\devnum{1},\ldots,\devnum{20}$
on all five development overlays; overlay~\devnum{0} alone continues through
$k=\devnum{76}$. This gives \devnum{156} capacity cells. Each replay starts at the
beginning of its complete trace, keeps queue state across weeks and drains all jobs;
there is no time compression, window reset or thinning. The same source jobs occur
in all overlays, so these are alternative constructions rather than independent
deployments. The grid and comparison family were declared after a pilot on the
first overlay, before the full sweep, not before inspecting the development demand.
% Sources: prechecks/weakness1_attack/summary_numbers.json, capacity ranges;
% prechecks/weakness1_attack/verification.json, complete_cells;
% prechecks/weakness1_attack/REPORT.md, declaration and state-continuity protocol.

Table~\ref{tab:s_capacity} reports the average of the overlay-specific deadline
p99 values, with pointwise intervals. Gap closure in this study is the ratio of
differences of those averaged quantiles. This differs from the package main tables,
which average overlay-specific gap ratios. Paired calendar-week resampling uses
\devnum{2{,}000} common draws, recomputing quantiles from the weighted job
distribution, not averaging weekly quantiles. The declared simultaneous band
covers \devnum{80} absolute FCFS-minus-policy differences: each capacity and each
non-oracle policy. It uses a centred coordinate-standardised maximum statistic.
These intervals resample realised outcomes, not queue trajectories; they omit
source-term selection, predictor fitting, shared source jobs, model mismatch,
feedback and cross-week state dependence. Maxima below are trace maxima, not
population confidence bounds.
% Sources: prechecks/weakness1_attack/summary_numbers.json, simultaneous_band;
% prechecks/weakness1_attack/capacity_curve.csv, p99_deadline_s and gap fields.

\begin{table}[htbp]
\centering\small
\caption{Capacity sweep with archived scores: mean overlay deadline p99 in seconds,
with pointwise paired calendar-week intervals. Negative effects are retained.}
\label{tab:s_capacity}
\begin{adjustwidth}{-\extralength}{0cm}
\centering
\begin{tabular}{rlll}
\toprule
$k$ & FCFS & SPJF-E & Guard(600) \\
\midrule
\devnum{1} & \devnum{108{,}345.31} [\devnum{92{,}267.38}, \devnum{114{,}125.36}] & \devnum{130{,}231.12} [\devnum{65{,}565.87}, \devnum{155{,}632.22}] & \devnum{108{,}839.10} [\devnum{92{,}739.31}, \devnum{114{,}582.34}] \\
% Source: prechecks/weakness1_attack/capacity_curve.csv, k=1, p99_deadline_s/lo/hi.
\devnum{2} & \devnum{10{,}362.62} [\devnum{7{,}565.14}, \devnum{11{,}503.45}] & \devnum{1{,}844.21} [\devnum{766.70}, \devnum{3{,}990.37}] & \devnum{10{,}618.06} [\devnum{7{,}820.32}, \devnum{11{,}755.23}] \\
% Source: prechecks/weakness1_attack/capacity_curve.csv, k=2, p99_deadline_s/lo/hi.
\devnum{3} & \devnum{1{,}056.40} [\devnum{688.91}, \devnum{1{,}273.01}] & \devnum{138.86} [\devnum{105.79}, \devnum{169.05}] & \devnum{934.68} [\devnum{452.24}, \devnum{1{,}253.73}] \\
% Source: prechecks/weakness1_attack/capacity_curve.csv, k=3, p99_deadline_s/lo/hi.
\devnum{4} & \devnum{273.50} [\devnum{197.09}, \devnum{352.10}] & \devnum{62.92} [\devnum{53.69}, \devnum{72.13}] & \devnum{89.61} [\devnum{65.94}, \devnum{114.14}] \\
% Source: prechecks/weakness1_attack/capacity_curve.csv, k=4, p99_deadline_s/lo/hi.
\devnum{5} & \devnum{119.35} [\devnum{92.91}, \devnum{143.43}] & \devnum{43.39} [\devnum{37.94}, \devnum{47.54}] & \devnum{47.44} [\devnum{40.46}, \devnum{53.39}] \\
% Source: prechecks/weakness1_attack/capacity_curve.csv, k=5, p99_deadline_s/lo/hi.
\devnum{8} & \devnum{24.43} [\devnum{18.46}, \devnum{29.58}] & \devnum{13.21} [\devnum{9.58}, \devnum{15.80}] & \devnum{13.36} [\devnum{9.67}, \devnum{16.05}] \\
% Source: prechecks/weakness1_attack/capacity_curve.csv, k=8, p99_deadline_s/lo/hi.
\devnum{14} & \devnum{0.84} [\devnum{0.45}, \devnum{1.50}] & \devnum{0.47} [\devnum{0.30}, \devnum{0.73}] & \devnum{0.47} [\devnum{0.30}, \devnum{0.73}] \\
% Source: prechecks/weakness1_attack/capacity_curve.csv, k=14, p99_deadline_s/lo/hi.
\devnum{20} & \devnum{0.00} [\devnum{0.00}, \devnum{0.04}] & \devnum{0.00} [\devnum{0.00}, \devnum{0.03}] & \devnum{0.00} [\devnum{0.00}, \devnum{0.03}] \\
% Source: prechecks/weakness1_attack/capacity_curve.csv, k=20, p99_deadline_s/lo/hi.
\bottomrule
\end{tabular}
\end{adjustwidth}
\end{table}


The negative cells are material. At $k=\devnum{1}$, FCFS-minus-Guard p99 differences
are \devnum{-208.61}, \devnum{-493.79} and \devnum{-1{,}071.72} s for the
\devnum{300}, \devnum{600} and \devnum{1200} s promises; all simultaneous bands
are wholly negative. At $k=\devnum{2}$ all guard point estimates remain negative.
Positive simultaneous lower bounds occur for SPJF-E at $k=\devnum{2}$--\devnum{12},
Guard(\devnum{300}) and Guard(\devnum{600}) at $k=\devnum{4}$--\devnum{12}, and
Guard(\devnum{1200}) at $k=\devnum{3}$--\devnum{12}. At $k=\devnum{14}$ the guard
point improvements are about \devnum{0.37} s but their simultaneous bands cross
zero; at $k=\devnum{20}$ the point differences are zero. A correct per-job bound
does not imply positive tail benefit across queueing capacities.
% Sources: prechecks/weakness1_attack/capacity_curve.csv, Guard rows,
% absolute_p99_improvement_s, simultaneous_improvement_lo/hi;
% prechecks/weakness1_attack/summary_numbers.json, support_by_policy.

For Guard(\devnum{600}) at $k=\devnum{4}$ the maximum excess is \devnum{515.84} s
and harm is \devnum{339.06} s, although it closes
\devnum{0.8000 [0.7620, 0.8388]} of the reference gap. At $k=\devnum{20}$ that
ratio is undefined because the reference gap is not positive, while maximum excess
and harm remain \devnum{124.38} and \devnum{37.37} s. On overlay~\devnum{0}, FCFS
still has one positive wait at $k=\devnum{75}$; at $k=\devnum{76}$ every policy has
zero wait. The latter is the exact maximum infinite-server occupancy, hence a
zero-wait endpoint for this overlay, not a capacity recommendation for a deployment.
% Sources: prechecks/weakness1_attack/capacity_curve.csv, Guard(600), k=4,20,
% max_excess_s, harm_s, gap_closed and interval; extension_curve.csv, k=75,76;
% prechecks/weakness1_attack/pilot.json and summary_numbers.json, endpoint.

\subsection{An Executing Timed-Work Service}

The physical selection rule chooses the complete calendar-aligned
\devnum{300}-second bin with the greatest offered capped work in development
overlay~\devnum{0}, breaking ties by earliest time. Its \devnum{2{,}151} jobs carry
\devnum{4{,}118.215007} s of requested work, \devnum{6.864} times the
two-worker release-window capacity. Each policy starts empty and drains its own
queue. Inter-arrival spacing, requested costs and scores are retained; workers hold
their slots with timed sleeps. This tests queueing and the control rule, not the
portability of submitted-program execution time. It also discards inherited
full-trace backlog, so its percentile is not the capacity sweep's percentile.
% Source: prechecks/weakness1_attack/physical_numbers.json, input.jobs,
% input.offered_requested_work_s/(input.window_s*input.parameters.workers), protocol.

The nominal payload cap is \devnum{60} s; the measured-service premise
$L=\devnum{61}$ s was fixed before execution. Guard parameters are
$B_0=\devnum{30}$ s, $\eta=\devnum{0.5}$, $\gamma=\devnum{0}$ and
$B_{\max}=\devnum{356}$ s, corresponding to the nominal
\devnum{300}-second ideal-model promise. A measured holding time above $L$ would
fail the premise, not be clipped or used to raise $L$ retrospectively. None did.
Actual enqueue, choice, dispatch, worker start/finish and completion receipt are
recorded separately; service becomes known only at completion receipt.
% Source: prechecks/weakness1_attack/physical_numbers.json, input.parameters.guard;
% prechecks/weakness1_attack/physical_verification.json, service-limit checks.

\begin{table}[htbp]
\centering\small
\caption{Physical timed-work runs. The ideal comparator uses each run's own measured
arrivals and holding costs; all entries are seconds.}
\label{tab:s_physical}
\begin{tabular}{lrrr}
\toprule
Policy & Mean wait & p99 wait & Max absolute per-job ideal discrepancy \\
\midrule
FCFS & \devnum{1{,}337.05} & \devnum{1{,}763.30} & \devnum{0.49} \\
% Source: prechecks/weakness1_attack/physical_numbers.json, policies.FCFS, wait_s and absolute_physical_minus_ideal_wait_s.
SPJF-E & \devnum{142.79} & \devnum{1{,}687.47} & \devnum{392.84} \\
% Source: prechecks/weakness1_attack/physical_numbers.json, policies.SPJF-E, wait_s and absolute_physical_minus_ideal_wait_s.
Guard(300) & \devnum{480.39} & \devnum{1{,}783.65} & \devnum{1{,}354.10} \\
% Source: prechecks/weakness1_attack/physical_numbers.json, policies.Guard(300), wait_s and absolute_physical_minus_ideal_wait_s.
\bottomrule
\end{tabular}
\end{table}


All \devnum{6{,}453} observed-prefix decisions passed an independent choice audit.
Physical and ideal measured-input Guard firing counts were \devnum{322} and
\devnum{372}. The small mean discrepancies do not establish job-level fidelity:
the maximum absolute per-job error is \devnum{392.84} s for SPJF-E and
\devnum{1{,}354.10} s for Guard. Guard's physical p99 exceeds the separate physical
FCFS run by \devnum{20.34} s. Thus the physical window confirms an executing
queue and rule conformance but does not show a guarded p99 gain. Actual arrivals
and holding costs differ between physical runs; their comparison is descriptive.
One selected window supplies no across-window confidence interval.
% Sources: prechecks/weakness1_attack/physical_numbers.json, policies.*.audit,
% absolute_physical_minus_ideal_wait_s and wait_s; physical_verification.json,
% Guard(300) ideal forced_count. Difference=1783.64837275-1763.3048828.

The guard's theorem reference is shadow FCFS on that guard run's measured arrivals
and service, not the separately timed physical FCFS run. An observed-trace
certificate augments the ideal bound by cumulative idle worker-time and
dispatch-to-start delay. Its largest correction is \devnum{1.127} s;
all \devnum{2{,}151} jobs pass, and the largest uncorrected excess over shadow FCFS
is \devnum{222.21} s. This is an ex post certificate conditional on the measured
trace and the \devnum{61} s service premise. It does not bound future operating-system
stalls and does not restore an advance uncorrected real-time promise.
% Sources: prechecks/weakness1_attack/stall_bound_check.json, metrics.
% maximum_additive_correction_twice_ns/2e9=1.1265093, corrected_violations=0;
% prechecks/weakness1_attack/physical_verification.json, Guard(300).max_exact_excess_s.

\subsection{Why Open-Loop Benefit Has No Known Bias Direction}

The replay compares policies conditional on fixed offered work. In an exact
single-server feedback construction, jobs of length \devnum{10} and \devnum{1} s
arrive together, the long job first, and the short job's user submits a second
unit job \devnum{0.5} s after completion. Replaying the FCFS-induced arrivals
gives SJF a mean-wait advantage of \devnum{3.00} s; letting each policy generate
its own successor arrival changes the advantage to \devnum{-0.167} s.
The corresponding linearly interpolated p99 advantages are \devnum{8.82} and
\devnum{0.47} s. The job cohort is identical, so this example does not depend on
comparing different populations.
% Source: prechecks/weakness1_attack/out_feedback_counterexamples.txt, EXAMPLE 1,
% exact Fraction schedules and open_loop/closed_loop mean and linear_p99 gains.

The opposite bias direction is possible too. An initial FCFS order with services
\devnum{20}, \devnum{1}, \devnum{3} and \devnum{10} s, whose unit-job user returns
after \devnum{2} s with a \devnum{0.1} s successor, gives open- and closed-loop
mean advantages of \devnum{9.20} and \devnum{11.16} s and p99 advantages of
\devnum{10.00} and \devnum{10.18} s. These exact recursions falsify a general
lower-bound interpretation; they are not fitted behavioural models. The pathwise
theorem still compares a guarded schedule with shadow FCFS on that schedule's
realised arrivals and services, not independently evolving closed-loop systems
\cite{shmueli2009simulation}.
% Source: prechecks/weakness1_attack/out_feedback_counterexamples.txt, EXAMPLE 2,
% parameters, exact schedules and gains; final assertion status PASS.

\subsection{LPC-EGEE: Recorded Queueing with Weak Counterfactual Validation}

The public LPC-EGEE trace supplies recorded waits and concentrated execution costs
together \cite{lpc_egee_trace}. After outage exclusion and queue-limit capping,
the longest \devnum{1\%} of \devnum{206{,}222} retained jobs carries
\devnum{40.7\%} of work. On the chronological evaluation part the top share is
\devnum{35.1\%}, which is the value used in the main cross-domain table.
The physical pools stay disjoint. Effective capacities \devnum{83} and
\devnum{56} are fitted only on the earlier part, with no test-driven refit.
Service is the minimum of recorded run time and its documented queue limit;
the common bound uses the largest pool limit, \devnum{259{,}200} s.
% Sources: prechecks/lpc_egee/applicability.csv, all/test aggregate rows;
% prechecks/lpc_egee/capacity.json and policy_parameters.csv.

The fitted FCFS approximation underestimates evaluation-part mean waits by
\devnum{24\%} and \devnum{39\%}, with hourly correlations of \devnum{0.44} and
\devnum{0.40}. The archive describes plain FCFS serial-job ordering during the
logged period, but queue concurrency ceilings, polling and middleware are absent
from the work-conserving replay. Recorded waits therefore cannot be equated with
the theoretical FCFS counterfactual. This is weak simulator validation, not
evidence that the original scheduler was an arbitrary non-FCFS policy.
% Source: prechecks/lpc_egee/validation.csv, fitted/test rows: mean ratios
% 0.757708/0.614029 and hourly Pearson 0.439849/0.397983; archive documentation.

In the conditional replay, SPJF-E reduces mean wait by
\devnum{21.4\% [13.2\%, 38.2\%]}, but its busy-hour p99 reduction of
\devnum{4.2\% [-11.2\%, 46.0\%]} remains unresolved. Busy hours use frozen
arrival-count thresholds, not deadlines. All six constant/age work-budget guards
at the tested multiples of $L$ fire zero times; a small fraction of the numerical
bound being spent is therefore not evidence of effective protection. Histories
use recorded completion times, and \devnum{2{,}906} evaluation arrivals
(\devnum{2.5\%}) have an outcome unavailable under replayed SPJF-E. This row is a
recorded-feature conditional counterfactual, not a closed-loop deployment result.
% Sources: prechecks/lpc_egee/pooled_policy_metrics.csv, SPJF-E mean and busy-p99
% reduction/interval fields; guard_dispatches.csv, all Constant/Age forced_sum=0;
% history_counterfactual_diagnostic.csv, affected count/share.


\providecommand{\devnum}[1]{#1}

\section{Additional Theory: Scope, Accounting and Instance-Dependent Bounds}
\label{sec:supp_theory_additions}

The results below delimit the main-text claims rather than enlarge them. The first
example rules out instance optimality for the uniform guard.
The exact identity that follows is a useful job-level accounting decomposition;
its core is standard sample-path workload balance, and we do not treat it as a
new conservation law. The final bounds replace the global timeout by realised
order statistics, but in their actual-size form they are retrospective
certificates under completion-only observation.

\subsection{Comparison with the Nearest Guarantee Disciplines}

\begin{table*}[htbp]
\centering
\caption{Scope of the nearest per-job guarantees and wrapper precedent.}
\label{tab:supp_related_comparison}
\scriptsize
\setlength{\tabcolsep}{2pt}
\begin{tabular}{p{0.11\textwidth}p{0.13\textwidth}p{0.10\textwidth}p{0.08\textwidth}p{0.15\textwidth}p{0.12\textwidth}p{0.18\textwidth}}
\hline
Method & Reference discipline & Arrival assumptions & Preemption & Information needed & Server model & Guarantee \\
\hline
This work & Same-input FCFS & Arbitrary finite path & No & Ranks, completions, public $L$, opaque base proposal & $k$ identical servers & Deterministic per-job additive waiting excess; wrapper around any base \\
PFCFS \cite{schwiegelshohn2000fairness} & Flow time without later submissions & Online releases & Gang preemption & Processor demand; service time need not be known; preemption penalty & Rigid parallel jobs & $\lambda$-fairness; PFCFS is $(2+2\bar s)$-fair with competitive efficiency bounds \\
Conservative backfilling \cite{lindsay2013backfilling} & Submission-time reservation & Online submissions & No & Requested wall time and resources & Rigid parallel jobs & Later backfills do not delay any reserved start \\
PV-EASY \cite{yuan2010pveasy,yuan2014strictfairness} & Priority reservation & Online submissions; prediction error allowed & Kill and restart & Runtime estimates and resources & Rigid parallel jobs & Zero lower-priority delay to the protected reservation \\
Nudge \cite{grosof2021nudge} & FCFS & Stationary Poisson & No & Size class at arrival & One server & Response-time stochastic dominance for its model; at most one swap per job \\
Load Balancing Guardrails \cite{grosof2019guardrails} & Unconstrained dispatcher & Stochastic heavy-traffic regime & SRPT at servers & Size class and dispatch state & Parallel server pools & Any dispatcher becomes mean-response-time optimal in heavy traffic \\
\hline
\end{tabular}
\end{table*}
% Source for the PFCFS definition and constants, and reading-depth notes for all
% rows: prechecks/guard_optimality_verify/literature_check.md and
% prechecks/guard_optimality_verify/verification.md, Section 9.

\subsection{No Uniform Instance-Optimality Ratio}

\begin{proposition}[An unbounded instance-optimality ratio]
\label{prop:supp_unbounded_instance}
At $k=1$, let three batch jobs have rank-order service times $(L,1,1)$, with
$L\geq2$, and set $G=L$. The cap of the parameter rule of the manuscript
(Corollary~3 there, $B_{\max} = k(G-(3-2/k)L)$) is $B=0$, so the
guard is FCFS and has total waiting time $2L+1$. The order $(1,1,L)$ has total
waiting time $3$ and satisfies the per-job promise $G$. Consequently the ratio
between the guard's total wait and that of a feasible policy is at least
$(2L+1)/3$, which is unbounded as $L\to\infty$.
\end{proposition}
% Source and direct checks at L=10,100,1000,10000:
% prechecks/guard_optimality_verify/out_t1.txt; derivation and scope in
% prechecks/guard_optimality_verify/verification.md, Section 4.

\begin{proof}
FCFS produces waits $0,L,L+1$. Under the competing order the two short jobs wait
$0$ and $1$, and the long job waits $2$. The long job's excess is therefore $2$
and both short jobs have nonpositive excess, so $L\geq2$ makes the order
$G$-feasible. The comparison allows a clairvoyant feasible policy; restricting
the comparator to completion-only policies would be a different question.
\end{proof}

The pointwise-maximal budget $B^{*}=G-L+1$ of Proposition~\ref{thm:k1_maximal}
admits one short job here, giving total wait $L+2$ and ratio $(L+2)/3$; the
conclusion is unchanged.

\subsection{Idle-Capacity Accounting}

Let $t_0$ be the first arrival time, let $b_Q(u)$ be the number of busy servers
under policy $Q$ at time $u$, and define the cumulative unused capacity
\[
 \Gamma_Q(t)=\int_{t_0}^{t}\bigl(k-b_Q(u)\bigr)\,du.
\]
This includes idle capacity caused by having fewer than $k$ jobs in the system;
work conservation excludes only deliberate idling while a job waits.

\begin{lemma}[Exact accounting with idle-capacity correction]
\label{lem:supp_gamma_identity}
For every finite input, every non-preemptive work-conserving policy $P$, and
every job $i$,
\[
 k\bigl(W_P[i]-W_{\mathrm{FCFS}}[i]\bigr)
 =\mathrm{In}_i-\mathrm{Out}_i
  +\Gamma_P(a_i)-\Gamma_{\mathrm{FCFS}}(a_i)
  -\rho_i^P+\rho_i^{\mathrm{FCFS}}.
\]
The identity was checked on \devnum{214{,}735} job-instances, including tied
events, non-batch arrivals, and same-phase overtakers.
\end{lemma}
% Source for the 214,735 checks and their decomposition:
% prechecks/guard_optimality_verify/out_t2.txt and
% prechecks/guard_optimality_verify/verification.md, Sections 7.1--7.2.

\begin{proof}
Let $A(t)$ be cumulative arrived work and $E_Q(t)$ cumulative executed work.
Then
\[
 U_Q(t)=A(t)-E_Q(t)
       =A(t)-k(t-t_0)+\Gamma_Q(t).
\]
At $a_i$, after simultaneous arrivals and before that phase's dispatches, work
from jobs not ranked below $i$ is identical in the two runs. Hence
\[
 R_i^P-R_i^{\mathrm{FCFS}}
 =U_P(a_i)-U_{\mathrm{FCFS}}(a_i)
 =\Gamma_P(a_i)-\Gamma_{\mathrm{FCFS}}(a_i).
\]
Substitution into the two busy-period equalities of the main text (the displayed equality $k\,W_P[i] = R^P_i + \mathrm{In}_i - \mathrm{Out}_i - \rho^P_i$ in the proof of the identity), one for $P$
and one for FCFS, proves the formula. A same-phase overtaker enters both
$\mathrm{In}_i$ and $\rho_i^P$ with its full size and therefore cancels.
\end{proof}

Writing $D_i=\Gamma_P(a_i)-\Gamma_{\mathrm{FCFS}}(a_i)$, the arrival-frozen
budget
\[
 b_i=\max\{0,\ k(G-L)-(k-1)L-D_i\}
\]
has the mathematical promise $G$ in the standard range
$G\geq(3-2/k)L$: combine
$\mathrm{In}_i<b_i+kL$ with Lemma~\ref{lem:supp_gamma_identity} and
$\rho_i^{\mathrm{FCFS}}\leq(k-1)L$. This is not a deployable rule under the
paper's completion-only observation model. In general
$\Gamma_{\mathrm{FCFS}}(a_i)$ depends on when an FCFS-shadow job completes
before its size has been revealed in the realised run; simulating the shadow
does not supply that missing completion time.

\subsection{Instance-Dependent Order-Statistic Bounds}

For a set of jobs $S$, let $\Lambda_r(S)$ be the sum of its largest at most $r$
original service times, with missing terms replaced by zero. We abbreviate
\[
 \Lambda_{k-1}(t)=\Lambda_{k-1}(\{j:a_j\leq t\}),\quad
 \Lambda^{<i}_{k-1}=\Lambda_{k-1}(\{j:j<i\}),
\]
and
\[
 \Lambda^{>i}_{k}(t)=\Lambda_k(\{j:j>i,\ a_j\leq t\}).
\]

\begin{proposition}[Instance-dependent workload gap]
\label{prop:supp_lambda_gap}
For any two non-preemptive work-conserving policies $Q_1,Q_2$,
\[
 |U_{Q_1}(t)-U_{Q_2}(t)|\leq\Lambda_{k-1}(t)
 \qquad\text{for every }t.
\]
\end{proposition}

\begin{proof}
The workload difference is continuous across arrivals, whereas
$\Lambda_{k-1}(t)$ can only increase. If
$U_{Q_1}(t)>\Lambda_{k-1}(t)$, policy $Q_1$ must have at least $k$ jobs present:
otherwise their remaining work is no larger than the sum of their at most
$k-1$ original sizes. All its servers are therefore busy, so the difference
cannot increase. It cannot cross the nondecreasing barrier from below; exchange
$Q_1$ and $Q_2$ for the other direction.
\end{proof}

\begin{proposition}[Instance-dependent excess accounting]
\label{prop:supp_lambda_excess}
For every policy $P$ and job $i$,
\[
 k\,\mathrm{excess}_P[i]
 \leq \mathrm{In}_i-\mathrm{Out}_i
 +\Lambda_{k-1}(a_i)+\Lambda^{<i}_{k-1}.
\]
\end{proposition}

\begin{proof}
In Lemma~\ref{lem:supp_gamma_identity}, Proposition~\ref{prop:supp_lambda_gap}
bounds the idle-capacity difference, equivalently the workload difference at
$a_i$, by $\Lambda_{k-1}(a_i)$. Drop $-\rho_i^P\leq0$. The at most $k-1$ FCFS
jobs still running when $i$ starts all have rank below $i$, so
$\rho_i^{\mathrm{FCFS}}\leq\Lambda^{<i}_{k-1}$.
\end{proof}

\begin{proposition}[Instance-dependent overtaking under the guard]
\label{prop:supp_lambda_in}
If job $i$ has an overtaker and $\tau$ is the dispatch time of its last
overtaker, then
\[
 \mathrm{In}_i
 <\mathrm{budget}(i,\tau)+\Lambda^{>i}_k(s_i^P).
\]
If the budget is nondecreasing in time, the first term may be replaced by
$\mathrm{budget}(i,s_i^P)$; with only a uniform cap it may be replaced by
$B_{\max}$.
\end{proposition}

\begin{proof}
Completed overtaking work before the last overtaker is dispatched is strictly
below $\mathrm{budget}(i,\tau)$. At $\tau$, at most $k$ dispatched overtakers
remain unfinished: the newly dispatched job and at most one on each other
server. Their total original size is at most
$\Lambda^{>i}_k(s_i^P)$.
\end{proof}

\begin{proposition}[Combined instance-dependent guard bound]
\label{prop:supp_lambda_guard}
Under a cap $B_{\max}$, every job satisfies
\[
 k\,\mathrm{excess}_P[i]
 \leq B_{\max}+\Lambda^{>i}_k(s_i^P)
 +\Lambda_{k-1}(a_i)+\Lambda^{<i}_{k-1},
\]
and the inequality is strict when $i$ has an overtaker. When it has none, the
sharper weak bound omits the first two terms.
\end{proposition}

\begin{proof}
Combine Propositions~\ref{prop:supp_lambda_excess} and
\ref{prop:supp_lambda_in}, use $\mathrm{Out}_i\geq0$, and apply the cap. If no
job overtakes $i$, then $\mathrm{In}_i=0$ and
Proposition~\ref{prop:supp_lambda_excess} gives the stated sharper form.
\end{proof}

These $\Lambda$ quantities use actual sizes, including sizes of unfinished
jobs, and hence are offline instance certificates in the completion-only model.
Replacing each unknown size by a public class-specific upper bound gives a
deployable but generally looser inequality. The homogeneous substitution recovers
the main text's aggregate worst-case constant; it does not prove that each
order-statistic term is separately minimal.

As an illustration, consider independently sampled truncated-Pareto sizes with
$\alpha=1.5$, $x_{\min}=1$, hard cap $1000$, sample size $N=200$, and
\devnum{5{,}000} samples drawn with seed \devnum{73621}. The first ratio below
sets $L$ after the fact to the maximum of each finite sample; the second uses the
fixed hard cap as the denominator.
% Source for all parameters and values in this paragraph and table:
% prechecks/guard_optimality_verify/out_t3.txt and
% prechecks/guard_optimality_verify/verification.md, Section 8.3.
\begin{center}
\begin{tabular}{rcc}
\hline
$k$ & $\mathbb{E}[\Lambda_{k-1}/((k-1)L_{\mathrm{sample}})]$
    & $\mathbb{E}[\Lambda_{k-1}/((k-1)L_{\mathrm{cap}})]$ \\
\hline
$4$  & $0.6825$ & $0.0444$ \\
$8$  & $0.4648$ & $0.0263$ \\
$16$ & $0.3139$ & $0.0161$ \\
\hline
\end{tabular}
\end{center}
Thus the larger ratios use $L=L_{\mathrm{sample}}=\max_j C_j$, not the published
timeout $L_{\mathrm{cap}}=1000$.

\subsection{Why Kill-and-Restart Tiering Does Not Escape the Floor}
\label{sec:supp_tiering}

Section~9 of the manuscript states that the obvious escape from the $\Omega(L)$ floor
fails. Kill-and-restart tiering runs every job for a short first quantum $\tau$ and
restarts the survivors, discarding the killed attempt; on $n$ jobs of length $L$ arriving
together it leaves some job with an excess of $n\tau/k - L$ over first-come first-served,
and that excess is wasted work rather than overtaking, so no budget removes it. The
light-job constant does not improve either, a second-tier piece being still a
non-preemptive block of length $L$: holding $\tau$ fixed and varying $L$, the worst
light-job additive constant we measure grows as about \devnum{1.1}$L$ and is flat in
$\tau$. The mechanism is TAGS \cite{harcholbalter2002tags} with the dedicated hosts
removed, and removing them is what breaks it: in TAGS the kill is a routing device and
the work is redone on the next host, whereas in a judging queue the kill is terminal.
% Ninth pass: this subsection stood in Section 9 of the manuscript, which keeps the
%   statement, the citation and the reason the mechanism breaks.

\subsection{Reservation-and-Refund Charging under Enforced Run-Time Caps}
\label{sec:supp_reservation}

Everything below assumes an information model the manuscript does not. Each job $j$
carries a cap $\ell_j$ with $C_j \le \ell_j \le L$ that is known when the job arrives and
that the service enforces, whereas Sections~5 and~6 of the manuscript assume a service
that learns a job's cost only at its completion. That is why this is a supplementary
subsection and not a
theorem of Section~6: admitting it into the main line would put two information models
side by side.

The mechanism itself is older than the guard, and we claim none of it. Conservative
backfilling gives every queued job a reservation when it is submitted and lets a job move
ahead only if it does not delay any job in the queue \cite{mualem2001backfilling}, and
enforced caps are standard there, a job that exceeds its declared run time being killed.
With one server-slot per job the head's reservation is whichever server is free at that
instant, so in our model both conservative and aggressive (EASY) backfilling reduce to
first-come first-served and neither is a comparator here. Lindsay
et~al.~\cite{lindsay2013backfilling} keep the same tentative profile, reserve on the
enforced estimate, compress the profile when a job ends early, and guarantee every job an
absolute start time fixed at its arrival. Reservation with refund is therefore their
mechanism; what differs here is the quantity guaranteed, a per-job bound on the delay
relative to first-come first-served, proved pathwise, with a maximal permission rule at
one server, and the exact constants at $k=1$. Load-balancing guardrails
\cite{grosof2019guardrails} charge dispatched work per size rank and per server and
balance those counters across servers, with the size known at dispatch, no reference
policy and no budget parameter; their Definition~2.1 does adjust a counter, lowering it
to the minimum across servers when a server empties, so it is not a rule without refund,
only one that never refunds at an individual job's completion.
% Source: prechecks/reservation_guard_verify/verification.md, Section 4 (the three
% readings, the corrected pointer to Definition 2.1, the counter reset, and the
% extension of the collapse argument to EASY) and its Section 5 list of statements that
% survive the check.

Write $\mathrm{over}_R[q](t)$ for the charge a waiting job $q$ carries: for every job of
rank above $q$ already dispatched, its cap while it is in service and its true service
time once it has completed.

\begin{proposition}[Reservation charging bounds the overtaken work]
\label{prop:supp_reservation_in}
Accept the base policy's proposal $j$ only when $j$ is the queue head or when
$\mathrm{over}_R[q]+\ell_j \le Z_q$ for every waiting $q$ of rank below $j$, and dispatch
the head otherwise. Then $\mathrm{In}_i \le Z_i$ for every input, every base policy and
every job $i$.
\end{proposition}

\begin{proof}
Rank is decided by $(a_j,\text{index})$, so every job counted in $\mathrm{In}_i$ is
dispatched after $i$ arrives and before $i$ starts, at an epoch where $i$ is waiting and
is not the head; the head exemption therefore does not apply and the test is run with
$q=i$. The charge carried for an already dispatched job is its cap while it runs and its
true service time afterwards, and both are at least its service time, so
$\mathrm{over}_R[i]$ at the dispatch of the $s$-th overtaker is at least
$\sum_{u<s}C_{j_u}$; the test then gives
$\sum_{u\le s}C_{j_u} \le \mathrm{over}_R[i]+\ell_{j_s} \le Z_i$. A rejected epoch
dispatches the head, whose rank is not above $i$, so it never enters $\mathrm{In}_i$. The
third step is where the cap has to be enforced: a job that may exceed $\ell_j$ breaks it
at once.
\end{proof}

\begin{proposition}[The promise under a constant allowance]
\label{prop:supp_reservation_promise}
With $Z \equiv B$, the rule gives $\mathrm{excess} \le B$ at $k=1$, attained, and
$k\,\mathrm{excess} < B+(2-2/k)L$ for $k \ge 2$, strictly. Setting
$B = k\bigl(G-(2-2/k)L\bigr)$ therefore offers the promise $G$.
\end{proposition}

The additive constant of the promise falls from $(3-2/k)L$ to $(2-2/k)L$, a drop of
exactly $L$ at every $k$: the $kL$ overshoot that completed-work charging pays in the
guard's overtaken-work bound disappears, while the $2(k-1)L$ slack of the identity
remains. That residual is still approached and not attained. On the witness family of
Section~S8.2 run at $B=f+L+1$, the ratio $(k\,\mathrm{excess}-B)/L$ is exactly
$2-2^{1-m}$ after $m$ rounds at $k=2$, reaching \devnum{255/128}, and reaches
\devnum{907/243} at $k=3$ and \devnum{323/64} at $k=4$, against the limits $2k-2$ of
\devnum{2}, \devnum{4} and \devnum{6}; a closed-form witness that attains the constant is
not known.
% Source: prechecks/reservation_guard_verify/verification.md, decision table rows C1-c,
% C1-d and C1-e, and out_c1b_indep.txt, which runs the S8.2 witness to m = 8.

Two consequences are worth separating, because they are easy to confuse. The
\emph{entry threshold} does not drop by $L$. An admitted overtaker always has
$\ell_j \le Z_q$, so at $k=1$ an overtaker of cap $L$ becomes admissible as soon as
$G \ge L$, which meets the $L$-scale floor of Proposition~1 of the manuscript with
equality, and for general $k$ the condition is $G \ge (2-1/k)L$ against
$G > (3-2/k)L$ for Algorithm~1 of the manuscript, a difference of $(1-1/k)L$, which is
zero at one server. It is the additive constant of the promise, not the threshold, that
falls by exactly $L$.

At one server the reservation is inert. A dispatch epoch there has an idle server, so no
job of rank above the head is in service and $\mathrm{over}_R$ coincides with the
completed-work charge: over \devnum{255{,}861} readings at $k=1$ the two never differ,
against \devnum{8{,}277} differences in \devnum{110{,}503} readings at $k \ge 2$. The
single-server gain comes entirely from the forward-looking $\ell_j$ term in the admission
test. For a constant allowance only the queue head has to be tested, because
$\mathrm{over}_R[q]$ is non-increasing in rank, so the implementation cost of
Section~5 of the manuscript is unchanged.
% Source: prechecks/reservation_guard_verify/out_overR_eq_overC_k1.txt (reading counts
% and mismatches at k = 1 and k >= 2) and out_c1b_indep.txt part 0 (19,381 head-only
% decisions, no divergence).

\begin{proposition}[One server: closed form and maximal permission]
\label{prop:supp_reservation_k1}
On the integer two-size batch family at one server, around an exact shortest-job-first
proposal sequence, the rule closes
$\min\bigl(1,\ \max(0,\lfloor (G-\ell)/\sigma \rfloor+1)/n\bigr)$ of the FCFS-to-SJF
total-wait gap. This is $\lfloor G/\sigma \rfloor/n$ at $\ell = C$, the size-aware
optimum, and $(\lfloor (G-L)/\sigma \rfloor+1)/n$ at $\ell = L$, the size-oblivious
frontier, so at $\ell=L$ it matches, and does not beat, the sharp completed-work guard.
Moreover, at any reachable history with queue head $h$, a wrapper that reads the caps and
promises $G$ must reject a later-ranked proposal $j$ whenever
$\mathrm{over}_R[h]+\ell_j > G$; the largest constant allowance with promise $G$ is
therefore exactly $B^{*}=G$, against $B^{*}=G-L+1$ for completed-work charging, and the
statement needs no integer lattice.
\end{proposition}

The maximality argument runs through the relabelling of $j$ to its cap. Three points
decide whether it closes. Only $j$ itself and the unfinished dispatched jobs of rank
above $h$ may be relabelled, and at $k=1$ there are none of the latter, so the general
form of the step is not available and is not needed. $W_{\mathrm{FCFS}}[h]$ is unchanged
because $h$ is the head, every job of rank below $h$ has already been dispatched, and
first-come first-served starts $h$ at a time fixed by the arrivals and service times of
jobs ranked below it alone. And the promise has to be quantified over the base policy's
proposal transcript: a base that reads true sizes, exact shortest-job-first among them,
could propose differently after the relabelling, and without that quantifier the
indistinguishability step fails. Nothing is needed about later steps, since $h$ is the
head, $\mathrm{Out}_h$ is always zero and $\mathrm{In}_h$ is non-decreasing.

At $k \ge 2$ the rule is not pointwise maximal, and this is settled rather than open.
Take two servers, four unit jobs arriving together and $G=0$. At the first epoch the rule
refuses every later-ranked proposal, because $\mathrm{over}_R[0]+\ell_1 = 1 > 0$, yet
dispatching job~1 is safe: the second server takes the head in the same instant, the
remaining two start at $t=1$ exactly as under first-come first-served, and every job has
excess $0$ under any continuation. The scale of the gap is not marginal. Of
\devnum{452{,}810} refused actions at $k=2$, \devnum{416{,}368} leave the victim inside
the promise after relabelling, and \devnum{481{,}964} of \devnum{491{,}340} do so at
$k=3$. This is what one expects once $\mathrm{In}$ and excess stop corresponding one to
one on several servers.
% Source: prechecks/reservation_guard_verify/verification.md, row C2-c and Section 2,
% counterexample 1; out_c2b_indep.txt, refused-action and safe-after-relabelling counts
% at k = 2 and k = 3.

Finally, the residual the reservation cannot remove is not an artefact of charging. There
are inputs on which a job has $\mathrm{In}_i=\mathrm{Out}_i=0$ and excess
$L\bigl(1-((k-1)/k)^m\bigr)$, so no charging rule of any kind removes an $\Omega(L)$
term. Its supremum is $L$, approached and not attained, which is the whole of
$(2-2/k)L$ at $k=2$ and only the part of it up to $L$ for $k \ge 3$.
% Source: prechecks/reservation_guard_verify/verification.md, rows C4a-a to C4a-c and
% out_c4b_indep.txt, k = 2 and k = 3 at m = 1..6 with excess/L = 1/2, 3/4, ..., 63/64
% and 1/3, 5/9, ..., 665/729.

The claims above are the ones an independent check re-derived from the written rule and
re-ran against its own implementation: \devnum{1{,}983{,}467} exhaustive runs and
\devnum{10{,}230{,}898} per-job checks give
$\max(\mathrm{In}-Z)=\devnum{0}$ with no violation, the instance form
$k\,\mathrm{excess} \le B+\Lambda_{k-1}(a_i)+\Lambda^{<i}_{k-1}$ holds with zero margin
over \devnum{1{,}488{,}984} checks, the closed form of
Proposition~\ref{prop:supp_reservation_k1} matches on all \devnum{848} family rows
including the corner cases $\ell = C$ and $\ell = L$, and the $k=1$ maximality step
survives \devnum{573{,}252} relabelling tests. Two statements of the earlier study did
not survive that check and are not made here: that the entry threshold drops by $L$, and
that maximality at $k \ge 2$ is open. A third needs a qualifier: the largest excess
attainable at budget $B$ equals $B$ when the cap is tight, $\ell = C$, and is $0$ while
$B < L$ when $\ell = L$, so in general only $\le B$ can be claimed.
% Source: prechecks/reservation_guard_verify/verification.md, decision table rows C1-a,
% C1-f, C2-a, C2-b, C3-d, C2-c and C6-c, with out_c1_indep.txt, out_t32_indep.txt,
% out_c2a_indep.txt and out_c2b_indep.txt.


%=================================================================
\begin{adjustwidth}{-\extralength}{0cm}
\reftitle{References}
% The same bibliography file as the manuscript; the class sets the style itself.
\bibliography{../docs/related_work/refs}
\end{adjustwidth}
\end{document}
